\section{Comparative Syntax and Language Paradigm Divergence}

After establishing Python's object model, frame mechanics, heap allocation, reference behavior, and CPython execution machinery, we can now return to the visible language surface. Syntax should not be treated as mere decoration. The way a language looks is often a compressed expression of how that language expects programs to be structured, executed, and understood.

This section compares Python's source-code shape with the C-style braced tradition. The purpose is not to rank the two syntaxes aesthetically, nor to repeat the shallow claim that Python is simply ``cleaner.'' The difference is deeper. C exposes more of the compiled, typed, storage-oriented model through declarations, braces, semicolons, pointers, and explicit memory-oriented operations. Python removes much of that visible machinery, but only because CPython carries more responsibility at runtime through objects, namespaces, frames, bytecode, dynamic type descriptors, and managed memory.

The comparison therefore begins at the surface, with code shape, but moves quickly underneath it. Braces versus indentation is not only a formatting preference. It reflects a different treatment of block structure by the lexer and parser. Python's lighter syntax is tied to runtime object behavior, while C's heavier syntax is tied to native compilation and fixed storage expectations. Understanding this divergence prevents the reader from translating Python mechanically into C-shaped ideas and helps explain why Python code behaves the way it does.